\documentclass{article}
\usepackage{amsmath,amssymb}
\usepackage{geometry}
\geometry{margin=1in}

\title{Step 224: Cross-Track CTMT Recursion Test on P vs NP}
\author{Six Birds Foundations III}
\date{2026-05-16}

\begin{document}
\maketitle

\section{Target}

This step tests whether the CTMT recursion pattern, already observed on RH, BSD, Hodge, and Navier--Stokes, also appears on the P-vs-NP complexity track.

\section{P-vs-NP Track State}

The P-vs-NP track is obstruction-oriented:
\[
\Xi_{\rm pack}>0
\]
is evidence that a declared package has not closed the SAT boundary.

The track is not at terminus. It is awaiting package specification for the selected next primary carrier
\[
\Pi_r^\#(\varphi)=\operatorname{lfp}(F_{\varphi,r}^\#).
\]
The sealed branch is a local-status/Tseitin proof-system branch with theorem-grade scoped obstruction. The active branch is the packaging-axis carrier
\[
P_{\Sigma,\tau}(3SAT).
\]
The active residual is
\[
\Xi_{\rm pack-atlas}(P),
\]
the gap between scoped packages and a lawful non-circular package atlas for polynomial-time visibility.

\section{Literature Audit}

Cook and Karp supply NP-completeness foundations. Baker--Gill--Solovay supply the relativization barrier. Razborov--Rudich supply the natural-proofs barrier. Aaronson--Wigderson supply the algebrization barrier. Williams supplies an algorithmic-diagonalization route for ACC lower bounds. Mulmuley--Sohoni supply the GCT program, while later results rule out some occurrence-obstruction routes. Proof-complexity lower bounds supply scoped leaves but do not cover all polynomial time without a bridge.

\section{Recursion}

The observed recursion is:
\[
\Xi_{\rm pack}
\to
\text{fixed quotient obstruction}
\to
\text{scoped local-status Tseitin branch}
\to
\text{proof-system bridge leaves}
\to
\Xi_{\rm atlas}(P)
\to
\text{universal P-atlas target-equivalence no-go}
\to
\text{packaging-axis level-1 host}
\to
\text{abstract-transformer package atlas}
\to
\text{no-smuggling/readability/atlas-scope gates}
\to
\text{barrier stack}.
\]

This is CTMT recursion in package-atlas terminal form. It is not literal matrix-element terminality.

\section{Verdict}

\[
\boxed{\texttt{V\_pvnp\_CTMT\_recursion\_verified}}
\]

The CTMT recursion pattern upgrades to \(\texttt{verified-on-5-track-instances}\) in structural form: RH, BSD, Hodge, NS, and P-vs-NP.

No P-vs-NP result is claimed.

\end{document}
